% GENERATED FILE. Edit canonical lessons, then run npm run generate:books.
% canonical-source-hash: fnv1a64:ef73f32d5d82bb14
% canonical-lessons: TA-C375-katci, TA-C375-kuttani, TA-C375-atikaram, TA-C375-celvakku, TA-C375-kavar

\chapter{Parties and Power}
\label{ch:ta-parties-and-power}
\hlchaptermodality{375}

\begin{chapteropening}
\textbf{By the end of this chapter:} \emph{I can talk about parties, alliances and influence.}

Complete the last lesson of chapter 375: I can talk about parties, alliances and influence.
\end{chapteropening}

\section[\texorpdfstring{\ta{கட்சி} (kaṭci)}{கட்சி (kaṭci)}]{\ta{கட்சி} (kaṭci) — a political party; a side}

\label{lesson:TA-C375-katci}



\begin{quote}
Before the new one: say the Tamil for a constituency, then the Tamil for a campaign.
\end{quote}

\subsection*{You'll want to know: \ta{கட்சி}}
\textbf{\ta{கட்சி}} — \emph{kaṭci} — \textquotedblleft{}a political party; a side\textquotedblright{}.

\begin{quote}
\ta{இந்தத்} \ta{தேர்தலில்} \ta{பல} \ta{கட்சிகள்} \ta{போட்டியிடுகின்றன}.
\end{quote}

\begin{quote}
\emph{intat tērtalil pala kaṭcikaḷ pōṭṭiyiṭukiṉṟaṉa.}
\end{quote}

\begin{quote}
Many parties are contesting this election.
\end{quote}

Plural: \ta{கட்சிகள்}, \emph{kaṭcikaḷ}.

Also a side in an argument: \ta{அவர்} \ta{என்} \ta{கட்சி}, he is on my side.

\begin{cousinweb}
\textbf{\ta{கட்சி}} (\emph{kaṭci}) is a word whose origin is not settled; it is often compared with Sanskrit \emph{pakṣa}, \textquotedblleft{}a wing, a side\textquotedblright{} (Monier-Williams), which also names a side in a debate.
\end{cousinweb}

\begin{grammarlens}[title={What you've built}]
One more word for giving reasons and opinions.
\end{grammarlens}

\subsection*{Your turn}
\begin{itemize}
\raggedright
  \item \emph{Say it:} \emph{kaṭci}
  \item \emph{Say it:} \emph{kaṭci}, once more
  \item \emph{Say it:} \emph{intat tērtalil pala kaṭcikaḷ pōṭṭiyiṭukiṉṟaṉa.} — Many parties are contesting this election.
  \item \emph{Recall:} say the Tamil for a constituency, then the Tamil for a campaign, then say \emph{kaṭci} again
\end{itemize}

\begin{tcolorbox}[breakable,colback=teal!4,colframe=teal!35!black,title={Before you move on}]
What does \ta{கட்சி} mean? (\textquotedblleft{}A political party; a side\textquotedblright{}.) Say it once more.
\end{tcolorbox}
\section[\texorpdfstring{\ta{கூட்டணி} (kūṭṭaṇi)}{கூட்டணி (kūṭṭaṇi)}]{\ta{கூட்டணி} (kūṭṭaṇi) — an alliance; a coalition}

\label{lesson:TA-C375-kuttani}



\begin{quote}
Before the new one: say the Tamil for a campaign, then the Tamil for a political party.
\end{quote}

\subsection*{You'll want to know: \ta{கூட்டணி}}
\textbf{\ta{கூட்டணி}} — \emph{kūṭṭaṇi} — \textquotedblleft{}an alliance; a coalition\textquotedblright{}.

\begin{quote}
\ta{இரண்டு} \ta{கட்சிகளும்} \ta{கூட்டணி} \ta{அமைத்தன}.
\end{quote}

\begin{quote}
\emph{iraṇṭu kaṭcikaḷum kūṭṭaṇi amaittaṉa.}
\end{quote}

\begin{quote}
The two parties formed an alliance.
\end{quote}

\ta{கூட்டணி} \ta{அமை}, to form an alliance.

\begin{cousinweb}
\textbf{\ta{கூட்டணி}} (\emph{kūṭṭaṇi}) is native: \emph{kūṭṭu}, \textquotedblleft{}union, joining\textquotedblright{} (DEDR 1882), and \emph{aṇi}, \textquotedblleft{}array of an army, row\textquotedblright{} (DEDR 117): ranks drawn up together.
\end{cousinweb}

\begin{grammarlens}[title={What you've built}]
One more word for giving reasons and opinions.
\end{grammarlens}

\subsection*{Your turn}
\begin{itemize}
\raggedright
  \item \emph{Say it:} \emph{kūṭṭaṇi}
  \item \emph{Say it:} \emph{kūṭṭaṇi}, once more
  \item \emph{Say it:} \emph{iraṇṭu kaṭcikaḷum kūṭṭaṇi amaittaṉa.} — The two parties formed an alliance.
  \item \emph{Recall:} say the Tamil for a campaign, then the Tamil for a political party, then say \emph{kūṭṭaṇi} again
\end{itemize}

\begin{tcolorbox}[breakable,colback=teal!4,colframe=teal!35!black,title={Before you move on}]
What does \ta{கூட்டணி} mean? (\textquotedblleft{}An alliance; a coalition\textquotedblright{}.) Say it once more.
\end{tcolorbox}
\section[\texorpdfstring{\ta{அதிகாரம்} (atikāram)}{அதிகாரம் (atikāram)}]{\ta{அதிகாரம்} (atikāram) — authority; power}

\label{lesson:TA-C375-atikaram}



\begin{quote}
Before the new one: say the Tamil for a political party, then the Tamil for an alliance.
\end{quote}

\subsection*{You'll want to know: \ta{அதிகாரம்}}
\textbf{\ta{அதிகாரம்}} — \emph{atikāram} — \textquotedblleft{}authority; power\textquotedblright{}.

\begin{quote}
\ta{இந்த} \ta{முடிவை} \ta{எடுக்க} \ta{அவருக்கு} \ta{அதிகாரம்} \ta{இல்லை}.
\end{quote}

\begin{quote}
\emph{inta muṭivai eṭukka avarukku atikāram illai.}
\end{quote}

\begin{quote}
He has no authority to make this decision.
\end{quote}

\ta{அதிகாரி}, an official, is the person who holds it.

\begin{cousinweb}
\textbf{\ta{அதிகாரம்}} (\emph{atikāram}) is borrowed from Sanskrit \emph{adhikāra}, \textquotedblleft{}authority; government, rule, office\textquotedblright{} (Monier-Williams), the source of \ta{அதிகாரி}.
\end{cousinweb}

\begin{grammarlens}[title={What you've built}]
One more word for giving reasons and opinions.
\end{grammarlens}

\subsection*{Your turn}
\begin{itemize}
\raggedright
  \item \emph{Say it:} \emph{atikāram}
  \item \emph{Say it:} \emph{atikāram}, once more
  \item \emph{Say it:} \emph{inta muṭivai eṭukka avarukku atikāram illai.} — He has no authority to make this decision.
  \item \emph{Recall:} say the Tamil for a political party, then the Tamil for an alliance, then say \emph{atikāram} again
\end{itemize}

\begin{tcolorbox}[breakable,colback=teal!4,colframe=teal!35!black,title={Before you move on}]
What does \ta{அதிகாரம்} mean? (\textquotedblleft{}Authority; power\textquotedblright{}.) Say it once more.
\end{tcolorbox}
\section[\texorpdfstring{\ta{செல்வாக்கு} (celvākku)}{செல்வாக்கு (celvākku)}]{\ta{செல்வாக்கு} (celvākku) — influence}

\label{lesson:TA-C375-celvakku}



\begin{quote}
Before the new one: say the Tamil for an alliance, then the Tamil for authority.
\end{quote}

\subsection*{You'll want to know: \ta{செல்வாக்கு}}
\textbf{\ta{செல்வாக்கு}} — \emph{celvākku} — \textquotedblleft{}influence\textquotedblright{}.

\begin{quote}
\ta{அவருக்கு} \ta{இந்த} \ta{ஊரில்} \ta{நிறையச்} \ta{செல்வாக்கு} \ta{இருக்கிறது}.
\end{quote}

\begin{quote}
\emph{avarukku inta ūril niṟaiyac celvākku irukkiṟatu.}
\end{quote}

\begin{quote}
He has a lot of influence in this town.
\end{quote}

X-\ta{இன்} \ta{மீது} \ta{செல்வாக்கு}, influence over X.

\begin{cousinweb}
\textbf{\ta{செல்வாக்கு}} (\emph{celvākku}) is usually read as native \emph{cel}, \textquotedblleft{}to go, pass, be current\textquotedblright{} (DEDR 2781), and \emph{vākku}, \textquotedblleft{}word\textquotedblright{}: a word that carries. The analysis is not settled.
\end{cousinweb}

\begin{grammarlens}[title={What you've built}]
One more word for giving reasons and opinions.
\end{grammarlens}

\subsection*{Your turn}
\begin{itemize}
\raggedright
  \item \emph{Say it:} \emph{celvākku}
  \item \emph{Say it:} \emph{celvākku}, once more
  \item \emph{Say it:} \emph{avarukku inta ūril niṟaiyac celvākku irukkiṟatu.} — He has a lot of influence in this town.
  \item \emph{Recall:} say the Tamil for an alliance, then the Tamil for authority, then say \emph{celvākku} again
\end{itemize}

\begin{tcolorbox}[breakable,colback=teal!4,colframe=teal!35!black,title={Before you move on}]
What does \ta{செல்வாக்கு} mean? (\textquotedblleft{}Influence\textquotedblright{}.) Say it once more.
\end{tcolorbox}
\section[\texorpdfstring{\ta{கவர்} (kavar)}{கவர் (kavar)}]{\ta{கவர்} (kavar) — to attract; to win over}

\label{lesson:TA-C375-kavar}



\begin{quote}
Before the new one: say the Tamil for authority, then the Tamil for influence.
\end{quote}

\subsection*{You'll want to know: \ta{கவர்}}
\textbf{\ta{கவர்}} — \emph{kavar} — \textquotedblleft{}to attract; to win over\textquotedblright{}.

\begin{quote}
\ta{அவருடைய} \ta{பேச்சு} \ta{பல} \ta{மாணவர்களைக்} \ta{கவர்ந்தது}.
\end{quote}

\begin{quote}
\emph{avaruṭaiya pēccu pala māṇavarkaḷaik kavarntatu.}
\end{quote}

\begin{quote}
His speech won over many students.
\end{quote}

Past (I …): \ta{கவர்ந்தேன்}, \emph{kavarntēṉ}.

Who is attracted takes -\ta{ஐ}. \ta{கவர்ச்சி}, attraction.

\begin{cousinweb}
\textbf{\ta{கவர்}} (\emph{kavar}) is native Dravidian: DEDR 1326 lists Tamil \emph{kavar}, \textquotedblleft{}to seize, grasp, get control of, desire\textquotedblright{}: to attract is to take hold of someone's mind.
\end{cousinweb}

\begin{grammarlens}[title={What you've built}]
One more word for giving reasons and opinions.
\end{grammarlens}

\subsection*{Your turn}
\begin{itemize}
\raggedright
  \item \emph{Say it:} \emph{kavar}
  \item \emph{Say it:} \emph{kavar}, once more
  \item \emph{Say it:} \emph{avaruṭaiya pēccu pala māṇavarkaḷaik kavarntatu.} — His speech won over many students.
  \item \emph{Recall:} say the Tamil for authority, then the Tamil for influence, then say \emph{kavar} again
\end{itemize}

\begin{tcolorbox}[breakable,colback=teal!4,colframe=teal!35!black,title={Before you move on}]
What does \ta{கவர்} mean? (\textquotedblleft{}To attract; to win over\textquotedblright{}.) Say it once more.
\end{tcolorbox}
